\Needspace{12\baselineskip}
\originalchapter{测量链、误差与可复现证据}
\originalsection{从物理量到读数}
实验比较的是带条件的物理模型和经过校准的观测。光子能量、电子动量、磁化强度通常不是仪器屏幕直接显示的量。装置把它们转成电荷、电压、计数或频谱，再用校准函数映回物理量。这条测量链中的每一环都可能引入增益、零点、带宽、阈值和非线性。

例如闪烁体将沉积能量转成光，光电倍增管把光转成电荷，前置放大器和成形放大器给出脉冲，计数器或多道分析器记录时间与幅值。只改一个放大增益，就能移动能谱峰的通道位置；物理峰能量本身并没有移动。先用已知谱线建立 $E=aC+b$，再测未知峰，才有可比较的能量值。

实验记录应在采集时保存原始读数、时间、装置几何、仪器型号、量程、阈值和校准版本。拟合脚本与结果要能回溯到这组输入。为说明方法而生成的模拟数据必须明确标为模拟，不可给它们编造实验日期、现场装置或真实样品身份。本册所有自行生成的数据都满足这一约定。

\originalsection{独立误差与共享误差}
设结果 $y=f(\mathbf x)$，输入协方差 $C$，局部线性误差传播为 $u_y^2=\sum_{ij}(\partial_i f)C_{ij}(\partial_j f)$。这是对小扰动的一阶近似；在接近零分母、阈值或强非线性区域，应传播完整概率分布或检查高阶项。相关误差不可通过多次重复自动消除。

若同一电压表零点漂移 $b$ 影响全部读数，$y_i=\theta+\epsilon_i+b$，独立项方差 $\sigma^2$、共享项方差 $u_b^2$，平均值方差是 $\sigma^2/N+u_b^2$。盲目把标准差除以 $\sqrt N$ 会把校准不确定度一起错误缩小。重复性良好也不能证明绝对准确。

直线测量 $y_i=a x_i+b$ 在独立已知 $u_i$ 条件下，用权重 $w_i=u_i^{-2}$ 最小化 $\chi^2=\sum_iw_i(y_i-a x_i-b)^2$。令 $X$ 的两列为 $x_i,1$，正规方程是 $X^TWX\hat\theta=X^TW\mathbf y$，参数协方差 $(X^TWX)^{-1}$。若 $x$ 也有显著误差，普通竖直残差拟合不再合适；共同校准量可以作为额外参数或协方差进入模型。

判断模型不能只看拟合曲线贴合。残差随 $x$ 的弯曲、随时间的漂移和相邻点相关，比一个漂亮的相关系数更能暴露错误。对计数数据直接用Poisson似然，低计数下把方差设为观测计数会使零计数点的权重失真。

\originalsection{装置使用和课程边界}
Junior Lab 的公开手册给出了研究型本科实验的原理、装置和分析要求，本册按测量问题改编本次实际取得的十九份实验手册。某些手册的课前网页题明确注明对OCW用户不可用；本册不捏造这些题。光源、磁场、放射源、高压和射频设备的具体现场操作应遵循负责实验室的规程；这里的示意图解释功能关系，不替代经认证的设备接线图。

一个可靠报告应说明目标量如何由可观察量获得、哪些参数独立校准、模型在哪个尺度有效以及数据是否支持结论。代码中的守恒量、解析极限与收敛比较提供相应计算证据，但数值通过也不代表实体实验已经做过。

\originalsection{一项实验如何形成可检验的问题}
实验的第一步是把目标分成可观察量与不可观察参数。例如检验光电关系时，可观察量是各滤光条件下的 $I(V)$ 曲线，截止电压是由曲线模型推断的量，$h/e$ 又是由多条截止电压对频率拟合得到的参数。若只留下最后六个截止读数而删除 $I(V)$，别人无法检验截止定义是否一致，也无法识别漏光造成的尾部。

预先写出模型、单位、独立标定及合理范围，能把“测到想要的答案”改成“数据可以怎样推翻模型”。对磁谱仪可同时比较 $p=mv$ 与 $p=\gamma mv$；对计数可比较Poisson、受死时间影响的过程和时间漂移。每个比较要使用同一组原始输入与误差模型，不能给不同理论任选最有利数据。

数据表中至少保存原始单位、采样时间、逐次读数、重复组、校准记录和删点理由。若模拟生成误差，必须连同分布、种子和真值参数保存。本卷代码把这些量放入参数和QA记录，课堂中无需实体仪器也能重现统计链，但这不替代手册规定的真实装置训练。

\originalsection{重复测量、标准误和自由度}
独立同分布读数 $x_i$ 的样本均值 $\bar x$ 和样本方差 $s^2=\sum(x_i-\bar x)^2/(N-1)$，分别估计位置和单次读数散布。标准误 $s/\sqrt N$ 衡量均值的不确定度，与仪器分辨率、单次读数误差是不同对象。小样本、正态模型且方差未知时，$(\bar x-\theta)/(s/\sqrt N)$ 服从Student $t_{N-1}$，区间不能总用固定的两倍标准误。

若序列相关，$\operatorname{Var}\bar x=\sigma^2[N+2\sum_{k=1}^{N-1}(N-k)\rho(k)]/N^2$。正相关减小有效样本数。将一秒内采一万个数字当一万次独立测量，可能严重低估噪声。分组、时间相关函数和不同采样间隔的比较能暴露这一问题。

系统误差不一定是未知的确定偏差。校准参数在给定先验或标定证据下也可用标准不确定度表示，并进入协方差；但没有根据的误差值不能靠“保守取大一点”伪装成统计结论。报告要给误差来源和计算依据。

\originalsection{直线拟合的解析结果和可识别性}
定义 $S=\sum w_i,S_x=\sum w_ix_i,S_y=\sum w_iy_i,S_{xx}=\sum w_ix_i^2,S_{xy}=\sum w_ix_iy_i$，$\Delta=SS_{xx}-S_x^2$。正规方程解为
\[
\hat a=\frac{SS_{xy}-S_xS_y}{\Delta},\qquad
\hat b=\frac{S_{xx}S_y-S_xS_{xy}}{\Delta},
\]
\[
u_a^2=S/\Delta,\quad u_b^2=S_{xx}/\Delta,\quad
\operatorname{Cov}(a,b)=-S_x/\Delta.
\]
当所有 $x$ 几乎相同，$\Delta\to0$，斜率不可识别；多次重复同一频率不能代替扩大频率覆盖。把横轴中心移到加权平均附近可以减少参数相关，虽不改变实际预测不确定度。

若每个观测方差是已知物理标准差，协方差无需另乘 $\chi^2/(N-2)$；若只知道相对权重、共同未知尺度从残差估计，才需要该因子。含相关误差时将 $W$ 换成 $C^{-1}$；若协方差近奇异，先检查是否把同一个共享参数重复计入。

残差 $r_i=y_i-\hat a x_i-\hat b$ 的方差受到拟合约束影响，不能把每个 $r_i/u_i$ 都当独立标准正态。$\chi^2$ 自由度为 $N-p$ 还依赖线性/局部线性模型及正确噪声假设。拟合图、残差图和误差来源应共同支持结论。

\originalsection{实验报告的论证链与计算文件}
报告正文按问题、装置功能、模型、数据、分析和结论连接；把仪器清单放进方法，但不要让它替代装置怎样把目标量转成读数的解释。图的坐标包含单位与误差；拟合线与数据点区分；结果给恰当有效数字。例如 $632.9\pm2.5$ nm已足够，写更多小数不提高信息。

计算工作流进一步记录输入哈希、软件版本、单位、每个输出的物理含义和运行命令。12.010 Fall2024的公开编程材料在本卷只用于这一工作流：数组、读写和脚本是工具，物理模型与适用范围仍是验收对象。没有把数学仓库已有的通用数值方法重新编成一门课程。
